\documentclass[11pt, a4paper]{article}

\usepackage{fontspec}
\usepackage{geometry}
\usepackage{fancyhdr}
\usepackage[hidelinks]{hyperref}
\usepackage[normalem]{ulem}
\usepackage{multicol}
\usepackage{enumitem}
\usepackage{amsmath}

\geometry{
  top=2cm,
  bottom=2cm,
  left=2cm,
  right=2cm,
  marginparsep=4pt,
  marginparwidth=1cm
}

\renewcommand{\headrulewidth}{0pt}
\pagestyle{fancyplain}
\fancyhf{}
\lfoot{}
\rfoot{}

\setlength{\parindent}{0pt}
\setlength{\parskip}{0pt}

\usepackage{xunicode}
\defaultfontfeatures{Mapping=tex-text}

\setlist[itemize]{leftmargin=1.4em, itemsep=0pt, topsep=1pt, parsep=0pt}

\newcommand{\code}[1]{\texttt{#1}}

\begin{document}

\small
\sloppy

\begin{center}
\normalsize\textbf{DSST389: Exam III --- Reference Sheet}
\end{center}

\smallskip

\textit{This sheet is attached to your exam. You do not need to memorize any of it.}

\smallskip

\textbf{Building a network.}
A network is a set of \textbf{nodes} and a set of \textbf{edges} joining pairs of them. It is
recorded as an \textbf{edge list}: one row per edge, with the endpoints in the columns
\code{doc\_id} and \code{doc\_id2}.

\begin{verbatim}
node, edge, G = DSNetwork.process(edges, directed=False)
\end{verbatim}

\code{node} has one row per node: \code{id}, the layout coordinates \code{x} and \code{y}, and
the metrics below.

\smallskip

\textbf{Components.}
A \textbf{component} is a set of nodes all reachable from one another by following edges. The
\code{component} column numbers them in \textbf{descending order of size}, so component 1 is
always the largest, and \code{component\_size} gives the size of each.

\smallskip

\textbf{Centrality.}
Distances, and the scores built on them, are computed \textbf{separately within each
component}. Filter to \code{component == 1} before comparing scores across nodes.

\begin{itemize}
  \item \code{degree} --- the number of neighbors a node has. Purely local.
  \item \code{close} --- \textit{closeness centrality.} Take the shortest-path distance from
        the node to every other node in its component, and add up the \textbf{reciprocals}. A
        neighbor contributes 1, a node four hops away contributes only $1/4$, so the score
        rewards being near \textit{everything}.
  \item \code{between} --- \textit{betweenness centrality.} For every pair of other nodes,
        consider the shortest paths between them; count what fraction of them run
        \textit{through} this node, and add that up over all pairs. High for
        \textbf{bridges}: nodes that are not popular but are unavoidable.
  \item \code{cluster} --- a community label, assigned by an algorithm that tries to put more
        edges inside groups than between them. It is an output, not a property: a node on the
        boundary between two groups may be assigned arbitrarily.
\end{itemize}

\smallskip

\textbf{Co-citation.}
An edge between two pages whenever some \textit{third} page cites \textbf{both}. Always
undirected. The \code{count} column holds the number of citing pages a pair shares.

\smallskip

\textbf{Images as arrays.}
\code{cv2.imread(path)} returns a NumPy array of shape \textbf{(height, width, channels)} with
values from 0 to 255. \code{img.mean()} averages over \textbf{every number}, channels included:
a $2 \times 2$ color image has 12 of them, not 4.

\begin{itemize}
  \item \textbf{OpenCV stores channels in BGR order, not RGB.} \code{[255, 0, 0]} is
        \textit{blue} and \code{[0, 0, 255]} is \textit{red}. \code{[0,0,0]} is black and
        \code{[255,255,255]} is white.
\end{itemize}

\smallskip

\textbf{Images as vectors.}

\begin{itemize}
  \item \code{SigLIPEmbedder()} --- embeds \textbf{images and text into the same space}, with
        \code{.embed\_image(path)} and \code{.embed\_text(string)}. The vectors are unit length,
        so \code{dot\_product(a, b)} is exactly the cosine similarity of Chapter 15: 1 means
        identical, 0 means orthogonal, and the angle distance is $\arccos$ of it.
  \item \textbf{Zero-shot classification}: embed each category description as text, and give
        each image the label whose text vector it has the largest dot product with. No training
        and no labeled examples --- an \textit{argmax over the candidate list you supplied}.
\end{itemize}

\end{document}
